% answers_short.tex —— 2009 年·普通高等学校招生考试·湖南卷·理科 · 简短答案速查
% 由 tools/make_short_answers.py 自 answers.tex 的 \gksolhead/\gkanswer 提取，
% 并与 2009/2009_数学真题_湖南理_答案要点.md 逐题比对后生成，勿手改（改源 MD 后重跑）。

\begin{gknotice}
\textbf{使用说明}：本节只给最终答案便于快速核对，完整推导见解析版。答案为编写者独立推导，\textbf{非官方标准答案}。
\end{gknotice}

\gkpart{选择题}
\begin{center}\begin{tabular}{ccccccccc}
\toprule
题号 & 1 & 2 & 3 & 4 & 5 & 6 & 7 & 8 \\
\midrule
答案 & D & A & D & B & C & B & C & D \\
\bottomrule
\end{tabular}\end{center}

\gkpart{填空题}
\begin{itemize}[leftmargin=2.6em]
  \item[\textbf{9.}] $12$。
  \item[\textbf{10.}] $8$。
  \item[\textbf{11.}] $2\sqrt2$。
  \item[\textbf{12.}] $\dfrac{\sqrt6}{2}$。
  \item[\textbf{13.}] $40$。
  \item[\textbf{14.}] （1）$12$；（2）$3$。
  \item[\textbf{15.}] $\dfrac{10}{3}$；$\dfrac{(n+1)(n+2)}6$。
\end{itemize}

\gkpart{解答题}
\begin{itemize}[leftmargin=2.6em]
  \item[\textbf{16.}] $A=30^\circ$，$B=120^\circ$，$C=30^\circ$，或 $B=30^\circ$，$C=120^\circ$。
  \item[\textbf{17.}] （1）$\dfrac16$；（2）分布列 $\dfrac1{27}$，$\dfrac29$，$\dfrac49$，$\dfrac8{27}$，$E\xi=2$。
  \item[\textbf{18.}] （1）证明见解析；（2）$\dfrac{\sqrt{10}}5$。
  \item[\textbf{19.}] （1）$y=m\sqrt x+\dfrac{256m}{x}+2m-256$（$x>0$）；（2）$9$ 个。
  \item[\textbf{20.}] （1）$\dfrac{x^2}{36}+\dfrac{y^2}{27}=1$（$x\geqslant2$），或 $y^2=12x$（$0\leqslant x<2$）；（2）$\dfrac{100}{11}$。
  \item[\textbf{21.}] （1）是；（2）②$\Rightarrow$④ 为真命题；（3）证明见解析。
\end{itemize}
